\section{Defining the Integers}
\begin{questions}
    % Defining the ``Shift Equivalence'' Relation
    \item Let $\EqRel$ be the relation on $\N^2$ where
    \[
        (a, b) \EqRel (c, d) \iff a + d = b + c
    \]
    for all $(a, b), (c, d) \in \N^2$. We show that $\EqRel$ is indeed an equivalence relation.
    \begin{itemize}
        \item \textbf{Reflexive}: Clearly $a + b = b + a$ for any $a, b \in \N$ (commutativity of addition, \myref{theorem:addition-commutative-naturals}), so $(a, b) \EqRel (a, b)$ by definition of $\EqRel$, proving reflexivity.
        
        \item \textbf{Symmetric}: Suppose $(a, b), (c, d) \in \N^2$ are such that $(a, b) \EqRel (c, d)$. This means $a + d = b + c$ by definition of $\EqRel$. Hence $c + b = d + a$ (commutativity of addition), so $(c, d) \EqRel (a, b)$, proving symmetry.
        
        \item \textbf{Transitive}: Suppose $(a, b), (c, d), (e, f) \in \N^2$ such that $(a, b) \EqRel (c, d)$ and $(c, d) \EqRel (e, f)$. This means
        \[
            a + d = b + c \;\text{ and }\; c + f = d + e.
        \]
        Adding $f$ to both sides of the first equation we get $a + d + f = b + c + f$. Note $b + c + f = b + d + e$ by the second equation, so we have
        \[
            a + d + f = b + d + e.
        \]
        Therefore $a + f = b + e$ by the cancellation property of addition (\myref{prop:addition-cancellation-property}), which means $(a, b) \EqRel (e, f)$, thus proving transitivity.
    \end{itemize}
    Thus $\EqRel$ is an equivalence relation.

    % Proving Partial Orders
    \item \begin{partquestions}{\alph*}
        \item \begin{proof}
            We are to prove that $\leq$ on $\Z$ is reflexive, antisymmetric, and transitive.
            \begin{itemize}
                \item \textbf{Reflexive}: Let $[(a,b)] \in \Z$. Clearly $a+b \leq a+b$ (since $\leq$ on $\N$ is a partial order and is thus reflexive), so $[(a,b)] \leq [(a,b)]$ by definition of $\leq$ on $\Z$.
                
                \item \textbf{Antisymmetric}: Suppose we have $[(a,b)], [(c,d)] \in \Z$ such that $[(a,b)] \leq [(c,d)]$ and $[(c,d)] \leq [(a,b)]$. This means
                \[
                    a + d \leq b + c \;\text{ and }\; c + b \leq d + a
                \]
                by definition of $\leq$ on $\Z$. But clearly the second relation is the same as $b+c \leq a+d$ (commutativity of addition), which thus means $a+d = b+c$ (as $\leq$ on $\N$ is antisymmetric). Therefore $[(a,b)] = [(c,d)]$ by definition of $\EqRel$ on $\Z$.
                
                \item \textbf{Transitive}: Suppose $[(a,b)], [(c,d)], [(e,f)] \in \Z$ are such that $[(a,b)] \leq [(c,d)]$ and $[(c,d)] \leq [(e,f)]$. Hence we have
                \[
                    a + d \leq b + c \;\text{ and }\; c + f \leq d + e.
                \]
                Adding $f$ to both sides of the first relation we get $a + d + f \leq b + c + f$. Note $b + c + f \leq b + d + e$ by the second equation, which thus means
                \[
                    a + d + f \leq b + d + e
                \]
                by transitivity of $\leq$ on $\N$. Therefore $a + f \leq b + e$, which means $[(a, b)] \leq [(e, f)]$.
            \end{itemize}
            Hence $\leq$ is indeed a partial order on $\Z$.
        \end{proof}

        \item \begin{proof}
            We will prove that $<$ is irreflexive and transitive, since asymmetry is implied by \myref{prop:asymmetric-iff-irreflexive-for-transitive}.
            \begin{itemize}
                \item \textbf{Irreflexive}: Let $[(a,b)] \in \Z$. Clearly $a + b \not< a + b$ (since $<$ on $\N$ is a strict partial order and is thus irreflexive), so $[(a, b)] \not< [(a, b)]$ by definition of $<$ on $\Z$.
                
                \item \textbf{Transitive}: Similar argument as in \textbf{(a)}, replacing $\leq$ with $<$.
            \end{itemize}
            Thus $<$ is a strict partial order on $\Z$.
        \end{proof}
    \end{partquestions}
\end{questions}
